\documentclass[twocolumn]{aastex701}

\usepackage{amsmath,amssymb}

\shorttitle{The modern parallax}
\shortauthors{Janzen}

\begin{document}

\title{The Modern Parallax: Can Matter Density Set the Rate of Cosmic Expansion?}

\author[gname=Daryl,sname=Janzen]{Daryl Janzen}
\affiliation{Department of Physics and Engineering Physics, University of Saskatchewan, 116 Science Place, Saskatoon, SK S7N 5E2, Canada}
\email[show]{daryl.janzen@usask.ca}

\begin{abstract}
Relativistic cosmology was built on the principle that matter determines how the universe expands. Stated globally, as in the Friedmann equation, the principle needs a cosmic time and spatial slices on which one density and one rate are defined, so it cannot test the structure it is stated on. Stated locally, as the averaging and timescape programmes take it, it needs no such structure and predicts that the expansion light records along each line of sight follows the matter it crosses. We test it with the cosmic microwave background (CMB), whose anisotropy is the sum of a source term fixed at last scattering and a cumulative term set by the expansion along each ray. If the expansion rate tracked the matter, with local rate $H(1-f\delta/3)$, the cumulative term on the Planck 2018 matter spectrum would scatter across the sky with an rms of $1.6\times10^{-2}$, some 390 times the whole measured anisotropy, and its angular power would exceed the measured spectrum at every multipole from $\ell=2$ to $2500$ by more than four orders of magnitude. Through its own equations, the timescape model exceeds it by 600 to 1300 times. Every simplification lowers the estimate, and the effect is distinct from the Sachs--Wolfe terms, which telescope where a differential expansion accumulates. The local density does not set the rate of cosmic expansion. With an observer-centred inhomogeneity closed by the Copernican principle, the cosmic time and the uniformity of its expansion, which the standard model posits before its dynamics, are measured.
\end{abstract}

\keywords{\uat{Cosmology}{343} --- \uat{Cosmic microwave background radiation}{322} --- \uat{Observational cosmology}{1146} --- \uat{Large-scale structure of the universe}{902} --- \uat{General relativity}{641}}

\section{Introduction}\label{sec:intro}

Can the density of matter shape the rate at which the universe expands? Relativistic cosmology was built on the answer yes. Einstein's first cosmological model fixed the size of the world from the density of its matter~\citep{Einstein1917}, and Friedmann's equations turned that relation into a law for the rate: the expansion is set by what the universe contains~\citep{Friedmann1922, Friedmann1924}. This is the reason a cosmology built on general relativity was expected to predict anything at all about the large-scale behaviour of the world. Whether the density sets the rate \emph{where the matter is}, as the local character of general relativity would lead one to expect, has nonetheless been settled by assumption rather than by measurement.

The assumption is built into the way the standard model is assembled (Section~\ref{sec:principle}). Before the density can set the rate there must be a rate and a density: a cosmic time whose slices are space, on which the expansion is one rate and the density one number. The standard model posits that structure first, as four kinematical assumptions, and states the dynamics last. In that global form the principle cannot test the structure, because it presupposes it, and every analysis that perturbs a Friedmann model inherits the same presupposition. The principle has one other form, and only one: a local form, in which the matter's own volume expansion responds to its own density along its worldlines. That form presupposes no slicing. It is the form the backreaction and timescape programmes take at its word (Section~\ref{sec:challenge}): a lumpy universe should have a lumpy expansion.

The local form makes a prediction that can be checked against something observed without first assuming a cosmic time. Light crossing a universe whose expansion follows the matter would accumulate different redshifts in different directions, according to the structure each ray crosses. The cosmic microwave background (CMB) carries that record. Its radiation reaches us across the whole of the expansion since decoupling, and the temperature received from each direction is the temperature at emission divided by the redshift accumulated on the way (Section~\ref{sec:decomp}). We compute the accumulated redshift the local form predicts, on the Planck 2018 matter spectrum, as a floor that every simplification lowers (Section~\ref{sec:floor}). It exceeds the measured anisotropy by more than two orders of magnitude in amplitude and four in angular power, and the timescape model, through its own equations, by several hundred to a thousand times. The answer to the question is no: in our universe the local matter density does not set the rate of cosmic expansion. With the one remaining escape, an inhomogeneity finely tuned about our position, closed by the Copernican principle (Section~\ref{sec:establishes}), the structure the standard model posits first is measured: a cosmic time and a uniform expansion on its slices. Section~\ref{sec:stands} sets out what the result leaves standing and what it removes, and Section~\ref{sec:follows} what follows from it.

This is the modern counterpart of a familiar episode. Before Bessel, whether the Earth moved was, on a strict reading, an interpretive choice among empirically equivalent pictures; stellar parallax converted it into a measured fact~\citep{Bessel1838}, and the heliocentric frame from an option into the one the measurement singled out. A purely relational reading of relativity holds, by the same token, that a preferred cosmic rest frame could never be observable even if it existed. The CMB overturns that expectation twice over. Its dipole measures our motion through the cosmic rest frame, and its isotropy, read as the record of the radiation's journey, measures the uniformity of the expansion in the cosmic time that frame carries.

We stress the register at the outset: nothing here modifies general relativity. The equations are unchanged. What is established is which of the formally available foliations is the physical one, and that the expansion it carries is not set locally by the matter.

\section{Density and the rate: how the principle entered cosmology}\label{sec:principle}

\subsection{From Einstein to Friedmann}\label{sec:history}

Einstein's cosmological paper of February 1917 rested on two assumptions, both drawn from experience~\citep{Einstein1917}. The first, an ``approximative assumption'' in his words, was that \emph{``there is a system of reference relatively to which matter may be looked upon as being permanently at rest''}---an assumption against the spirit of relativity~\citep{JanzenEinsteinCosmology}. His ground was empirical: ``the most important fact that we draw from experience as to the distribution of matter is that the relative velocities of the stars are very small as compared with the velocity of light''. The second was that the world is static. A static, finite world filled with matter collapses under its own gravity, and Einstein added the cosmological term $\lambda$ to the field equations to hold it in balance. The balance fixes the size of the world from what it contains: with $\lambda=\kappa\rho/2$ the radius is $R=1/\sqrt{\lambda}$, so \emph{the density of the world-matter determines the size of the world}. That is the first statement of the principle this paper tests, and it runs in one direction, from the content to the geometry. Eddington later showed that the balance is unstable~\citep{Eddington1930}.

De~Sitter's reply came one month later and countered every consideration in turn~\citep{deSitter1917}. His own solution of the field equations with $\lambda$ contained no matter at all. The boundary conditions were dismissed as unobservable in principle and, given the new $\lambda$-term, achievable anyway; the world-matter Einstein thought essential, de~Sitter wrote, ``serves no other purpose than to enable us to suppose it not to exist''; and the choice among the candidate universes was placed outside physics altogether: ``The choice can thus not be decided by physical arguments, but must depend on metaphysical or philosophical considerations, in which of course also personal judgment or predilections will have some influence''. The criticism he pressed most was the one that matters here: that in Einstein's solution ``time has a separate position'' from space. The matter's rest frame singles out a cosmic time.

Three years later the same structure was defended on geometric grounds by Eddington, who held that ``the world taken as a whole has one direction in which it is not curved; that direction gives a kind of absolute time distinct from space'', that relativity ``is not concerned to deny the possibility of an absolute time, but to deny that it is concerned in any experimental knowledge yet found'', and that ``it need not perturb us if the conception of absolute time turns up in a new form in a theory of phenomena on a cosmical scale, as to which no experimental knowledge is yet available''~\citep{Eddington1920}.

Friedmann dropped the requirement that the world be static~\citep{Friedmann1922, Friedmann1924}. Under the same assumptions of a cosmic time and a uniform space, the field equations then admit worlds that expand or contract, at a rate set by the density of their matter, the curvature of their space and $\lambda$. This is the form the principle has held ever since: Einstein's direction of determination, from content to geometry, carried over to the rate. Lema\^itre derived the expanding solutions independently and connected them with the recession of the nebulae~\citep{Lemaitre1927}, and the redshift--distance relation~\citep{Hubble1929} made expansion the reading of the data.

Einstein's own engagement with the question then fell nearly silent. The 1917 paper was the last in which he published the details of his cosmological considerations; after Hubble's result he wrote some five pages on the problem in two short articles, the second with de~Sitter, and his later rejection of the cosmological term shaped the expansion paradigm for the rest of the century~\citep{JanzenEinsteinCosmology}. He never returned in print to the question his own assumption raised: how a cosmic time distinguished by the matter's rest frame stands to the relativity of simultaneity.

\subsection{The order of construction}\label{sec:order}

Robertson and Walker gave the general class of spacetimes in which these models live~\citep{Robertson1929, Robertson1933, Robertson1935, Walker1937}, and the standard model of cosmology is assembled from it in a definite order~\citep{north1965}. One first posits that the average rest frame of matter---made manifest in the redshift--distance relation---defines a coherent cosmic time~\citep{Weyl1923}, the one direction in which the four-dimensional world is, in Eddington's phrase, not curved; then that space is the family of hypersurfaces orthogonal to it; then that this space is isotropic; and then that it is homogeneous. Only last does the geometry evolve, by a scale factor $a(t)$ whose rate the field equations fix through the energy content,
\begin{equation}
H^{2}\equiv\left(\frac{\dot a}{a}\right)^{2} = \frac{8\pi G}{3}\,\bar\rho - \frac{k}{a^{2}} + \frac{\Lambda}{3}.
\label{eq:friedmann}
\end{equation}
The first four steps are kinematical posits; the dynamics comes fifth and last. The standard presentation treats the kinematical posits as idealizations vindicated after the fact by the success of the fit---a global cosmic time, a preferred rest frame and a uniform expansion accepted because the model built on them works.

The order is not incidental, and it is the crux of what follows. Equation~\eqref{eq:friedmann} relates \emph{the} density to \emph{the} rate, and neither is defined until the first four steps have done their work. The rate is the rate at which one family of slices advances, so it needs the cosmic time and its slices. The density that enters is the density on a slice, and it is a single number only because isotropy and homogeneity make it the same everywhere on the slice. The principle that density sets the rate is therefore stated on a structure that must be posited before it can be stated. Nothing can be said about space, its density or its rate of expansion until a cosmic time and its slices have been assumed. In this global form the principle cannot test the kinematical posits, because it presupposes them. The same holds of every analysis that perturbs a Friedmann model, from cosmological perturbation theory to the stability arguments of Section~\ref{sec:tension}: each puts the dynamics to the test on a foundation already assumed.

\subsection{The two forms of the principle}\label{sec:twoforms}

There is one form of the principle that presupposes no slicing, and only one. Along the matter's own worldlines, with four-velocity $u^{a}$, the volume expansion $\Theta=\nabla_{a}u^{a}$ of the matter flow obeys the Raychaudhuri equation~\citep{Raychaudhuri1955},
\begin{equation}
\dot\Theta = -\tfrac13\Theta^{2} - \sigma_{ab}\sigma^{ab} + \omega_{ab}\omega^{ab} - 4\pi G(\rho+3p) + \Lambda + \nabla_{a}A^{a},
\label{eq:raychaudhuri}
\end{equation}
with $\sigma_{ab}$ the shear, $\omega_{ab}$ the vorticity and $A^{a}$ the acceleration of the flow, the dot denoting the derivative along $u^{a}$, and $c=1$. The local density enters directly. No global time, no slices and no averaging are needed to state it: the matter's expansion responds to the matter's density where the matter is. For a homogeneous and isotropic flow it reduces to the acceleration equation of the Friedmann models, of which~\eqref{eq:friedmann} is the first integral, so the global form is the local one restricted to the slices of the first four posits. The local form is the principle read as a local theory of gravity reads its field equations.

The local form makes a definite prediction for light. The redshift measured by the members of a matter congruence along a ray with spatial direction $e^{a}$ accumulates per unit proper distance $\ell$ as~\citep{Ellis1971, EllisMaartensMacCallum2012}
\begin{equation}
\frac{d\ln(1+z)}{d\ell} = \tfrac13\,\Theta + \sigma_{ab}e^{a}e^{b} + A_{a}e^{a}.
\label{eq:redshiftrate}
\end{equation}
In a Friedmann model only the first term survives and $\Theta=3H$. If the density sets the rate locally, in the sense that each region expands as a Friedmann world of its own density, then each region's expansion is isotropic, its shear vanishes and, for pressureless matter, its acceleration vanishes. The redshift then accumulates the matter's local expansion along the path, $d\ln(1+z)=\tfrac13\Theta\,dt=H_{\mathrm{loc}}\,dt$: the light records the expansion of the matter it crosses.

So the principle has exactly two forms. The global form needs the kinematical posits before it can be written down. The local form needs no slicing, and it predicts that the expansion light records follows the matter. Only the local form can be put to a test that does not first assume what is being tested, and this paper tests it.

\section{The tension and the challenge}\label{sec:challenge}

\subsection{Locality and a lumpy universe}\label{sec:tension}

General relativity is a local theory: spacetime tells matter how to move, and matter tells spacetime how to curve~\citep{Wheeler1998}, where the matter is. Read with local field equations, Einstein's and Friedmann's principle is local too: the matter sets the expansion where the matter is, and the global rate is its average. The standard account takes that reading implicitly whenever it speaks of matter decelerating the expansion.

The universe is not uniform. Over 13.8 billion years it has grown lumpy, with density contrasts of order unity and larger below about $10\,h^{-1}$\,Mpc and voids filling much of the volume, and it approaches statistical homogeneity only on scales of order $100\,h^{-1}$\,Mpc~\citep{Scrimgeour2012}. The standard justification for the kinematical posits is that on large enough scales the matter is distributed, and evolves, uniformly enough that the averages behave: the posits hold for the large-scale average, and the lumpiness is perturbation upon it. Cosmological perturbation theory carries this out. In it the matter's local expansion does vary with its density, but it varies as a perturbation about a background $a(t)$ given in advance. Green and Wald have shown that small-scale inhomogeneities backreact negligibly on the large-scale metric~\citep{GreenWald2011, GreenWald2014}. Their argument takes the metric to be close to Friedmann--Lema\^itre--Robertson--Walker form---a cosmic time and a single global expansion---and shows that the lumpiness does not disturb it. It establishes the \emph{stability} of the foliation it assumes, not its existence, and by Section~\ref{sec:order} it could not do otherwise: an analysis that perturbs a Friedmann model has already assumed the first four posits.

\subsection{Backreaction and timescape}\label{sec:backreaction}

The gap has been pressed directly. Ellis posed the fitting problem: which smooth model is the right one to fit to a lumpy universe~\citep{Ellis1984}? Buchert showed that in general relativity averaging and time evolution do not commute: the averaged expansion of an inhomogeneous dust universe obeys Friedmann-like equations with an additional \emph{kinematical backreaction} term $\mathcal{Q}$, set by the variance of the local expansion rates about their average~\citep{Buchert2000, Buchert2008}. R\"as\"anen, Kolb, Matarrese and Riotto, and others argued that this term might account for the apparent acceleration of the expansion without dark energy~\citep{Rasanen2004, Kolb2006}, and the averaging, backreaction and fitting problems have been reviewed at length~\citep{Clarkson2011, BuchertRasanen2012}. Whether the effect is small remains disputed~\citep{GreenWald2011, GreenWald2014, Buchert2015}. For light, R\"as\"anen showed that in statistically homogeneous dust spacetimes the redshift along a long ray tracks the volume-averaged expansion, the local departures of the rate from the average entering as a fluctuation argued to be small when integrated over long distances~\citep{Rasanen2009}.

The timescape model takes the programme furthest~\citep{Wiltshire2007, Wiltshire2007PRL, Wiltshire2009}. It divides the universe into voids and walls, with void volume fraction $f_{\rm v}$, whose expansion rates $H_{\rm v}$ and $H_{\rm w}$, referred to the volume-average time $t$, differ: the volume-average rate is $\bar H=\bar\gamma_{\rm w}H_{\rm w}=\bar\gamma_{\rm v}H_{\rm v}$, with $\bar\gamma_{\rm w,v}=dt/d\tau_{\rm w,v}$ the ratio of volume-average time to the proper time of each region. In its tracker solution the walls expand at two-thirds the rate of the voids, and the difference is the backreaction that produces its apparent acceleration, $\mathcal{Q}=6f_{\rm v}(1-f_{\rm v})(H_{\rm v}-H_{\rm w})^{2}$. Its further postulate is that each region measures the same expansion rate in its own proper time, the difference being carried by the rates of clocks, which differ between walls and voids by 38 per cent today. Supernova data have recently been argued to prefer it to the standard model~\citep{Seifert2025}.

What these programmes share is the local form of the principle, taken at its word: a lumpy universe should have a lumpy expansion. They take that position on the grounds of general relativity's own locality, against a standard model that applies the principle only to the average. The question of the introduction is therefore live, and it is empirical. One side holds that the uniform foliation is a large-scale idealization that works; the other that the local physics should make the expansion follow the matter. Only the second states the principle in a form that does not already assume the answer.

\subsection{What would decide it}\label{sec:decide}

To decide between them, the consequence of the local form must be worked out for something observed without first assuming a cosmic time. The nearest argument in the literature is of a different kind. The Ehlers--Geren--Sachs theorem and its almost-isotropic extension deduce the geometry from a symmetry: if every comoving observer sees an (almost) isotropic radiation background, the spacetime is (almost) Friedmann~\citep{EGS1968, Stoeger1995}. Its premise is isotropy at every observer, which cannot be observed from here, and the almost-isotropic form needs bounds on the derivatives of the anisotropies as well, without which an almost isotropic temperature does not imply an almost isotropic universe~\citep{Nilsson1999}. Neither reads the isotropy we do observe as what it also is: the record of the radiation's journey to us. The isotropy of the CMB has long been used to bound a global anisotropy of the expansion, a shear of the Bianchi type~\citep{CollinsHawking1973, Saadeh2016}, and direction-dependent expansion rates are now mapped directly, from the CMB and from low-redshift distances~\citep{Zhou2025}. The local form predicts the other way the expansion could fail to be uniform: a rate that varies from region to region with the density. The observable that decides it is the redshift of the CMB.

\section{The observable: the redshift as the record of a journey}\label{sec:decomp}

The cosmological redshift along a line of sight is the expansion the photon accumulated in transit. In general it is~\eqref{eq:redshiftrate} integrated along the ray; in a Friedmann model it is $1+z = a_{\mathrm{obs}}/a_{\mathrm{em}}$, equivalently $\ln(1+z) = \int H\,dt$ along the ray. Either way the redshift \emph{is} the accumulated expansion in that direction, a path integral and not a property of the source. In the Friedmann case the path integral is a conserved quantity written out. In conformal form $ds^{2}=a^{2}(\eta)(-d\eta^{2}+d\mathbf{x}^{2})$ the vector $\partial_\eta$ is a conformal Killing vector, so its charge $\xi\!\cdot\!p$ is conserved along null geodesics; for a photon that charge is $a\omega$, and its constancy is $1+z=a_{\mathrm{obs}}/a_{\mathrm{em}}$, whose logarithm is $\int H\,dt$.

The observed temperature is $T_{\mathrm{obs}}(\hat n) = T_{\mathrm{em}}(\hat n)/(1+z(\hat n))$, so the observed anisotropy separates exactly into two physically distinct contributions,
\begin{equation}
\frac{\delta T_{\mathrm{obs}}}{T_{\mathrm{obs}}} = \frac{\delta T_{\mathrm{em}}}{T_{\mathrm{em}}} \;-\; \delta\ln(1+z),
\label{eq:decomp}
\end{equation}
a \emph{source} term, the Sachs--Wolfe and acoustic anisotropies fixed by conditions on the last-scattering surface, and a \emph{cumulative} term, the directional variation of the expansion accumulated over the whole path. The two are independent: the value of one constrains the other not at all. Observed isotropy therefore requires \emph{both} to be small---uniformity at decoupling \emph{and} uniformity of the subsequent expansion. The standard reading attributes the observed anisotropy to the source term alone, the cumulative term having been set to zero in advance by the posit of a single scale factor, so it is never put to the test.

The decisive consequence is seen by granting that reading its strongest form. Suppose the last-scattering surface were perfectly homogeneous, $\delta T_{\mathrm{em}} = 0$ exactly. Then \eqref{eq:decomp} reads $\delta T_{\mathrm{obs}}/T_{\mathrm{obs}} = -\,\delta\ln(1+z)$, with nothing left over: the entire observed isotropy is the statement that the integrated expansion was the same in every direction. Homogeneity at the source does not \emph{supply} the isotropy; granted perfectly, it removes one term and leaves the observed isotropy as a \emph{direct measurement of uniform expansion}. The surface is not in fact homogeneous, and that sharpens the point: its Sachs--Wolfe and acoustic anisotropies, at the $10^{-5}$ level, arrive with their predicted spectrum intact. The journey that carried them imposed a redshift of 1089, and it imposed it the same in every direction to well within one part in $10^{5}$; otherwise it would have overwritten the primordial pattern. The common supposition that the isotropy of the CMB is delivered by homogeneity at last scattering credits the source alone for a fact about the \emph{observed} radiation that resides in both terms, the second of which only the intervening expansion can supply. The next section puts a number to the cumulative term under the local form of the principle, and reads that number against the data.

\section{The test: if the density set the rate, what would the sky show?}\label{sec:floor}

\subsection{The hypothesis}

The hypothesis is the local form of Section~\ref{sec:twoforms}: the expansion rate tracks the matter region by region, each region expanding as a Friedmann world of its own density, with no single global scale factor. This is the premise of the averaging programme, in which domains of different density expand at different rates~\citep{Buchert2000}, and of treatments of light propagation in which the redshift accumulates the local expansion rate along the ray~\citep{Rasanen2009}. By~\eqref{eq:redshiftrate}, the shear and acceleration of each region vanish and the redshift accumulates the local rate, $d\ln(1+z)=H_{\mathrm{loc}}\,dt$.

To linear order the local rate of a region of smoothed matter contrast $\delta_R$ follows from mass conservation, $\dot\delta=-\delta\Theta$, which gives
\begin{equation}
H_{\mathrm{loc}} = H\left(1-\tfrac13 f\,\delta_R\right),\qquad f\equiv\frac{d\ln D}{d\ln a},
\label{eq:localrate}
\end{equation}
with $D$ the linear growth factor. This is the separate-universe correspondence~\citep{Peebles1980}: an overdense region expands more slowly, an underdense one faster.

\subsection{Accumulation along the line of sight}

A photon crossing comoving distance $d\chi$ spends $dt=a\,d\chi$ in transit, so the direction-dependent part of its accumulated redshift is
\begin{equation}
\delta\ln(1+z)(\hat n) = -\int_0^{\chi_*} d\chi\; g(\chi)\,\delta_R\big(\chi\hat n,\,z(\chi)\big),
\label{eq:lineint}
\end{equation}
with $g(\chi)=\tfrac13\,aHf$, integrated to the last-scattering distance $\chi_*$. Two statistics of this field follow directly from the linear matter power spectrum $P_R(k,z)=P(k,z)\,W^{2}(kR)$, with $W$ the top-hat window of radius $R$. The variance along a single ray is
\begin{equation}
\sigma^{2}_{\mathrm{path}} = \int_0^{\chi_*} d\chi\; g^{2}(\chi)\,\frac{1}{2\pi}\int_0^\infty dk\,k\,P_R(k,z),
\label{eq:sigmapath}
\end{equation}
and its angular power spectrum is, in the Limber approximation~\citep{Limber1953},
\begin{equation}
C_\ell = \int_0^{\chi_*} d\chi\;\frac{g^{2}(\chi)}{\chi^{2}}\,P_R\!\left(\frac{\ell+1/2}{\chi},\,z(\chi)\right),
\label{eq:limber}
\end{equation}
which is directly comparable, multipole by multipole, with the measured temperature spectrum in units of $(\Delta T/T)^2$.

\subsection{The number}

We evaluate \eqref{eq:sigmapath} and \eqref{eq:limber} on the linear matter spectrum of the Planck 2018 base-$\Lambda$CDM model~\citep{Planck2018VI}, computed with \textsc{camb}~\citep{Lewis2000} ($\sigma_8=0.811$, $z_*=1089.9$), with the growth rate $f$ taken from the same spectrum. For a smoothing radius $R=8\,h^{-1}$Mpc,
\begin{equation}
\sigma_{\mathrm{path}} \simeq 1.6\times 10^{-2}.
\label{eq:value}
\end{equation}
The rms of the measured CMB temperature anisotropy, summed over $\ell=2$--$2508$ from the Planck 2018 spectrum~\citep{Planck2018V}, is $4.1\times10^{-5}$. The hypothesis predicts a scatter in the accumulated redshift alone some $390$ times larger than the entire measured anisotropy. Figure~\ref{fig:floor} makes the comparison scale by scale: the angular power \eqref{eq:limber} exceeds the measured power at every multipole from $\ell=2$ to $2500$, by a factor of at least $2\times10^{4}$ and of $2\times10^{5}$ at the median multipole.

\begin{figure}[t]
\centering
\includegraphics[width=\columnwidth]{fig_floor.pdf}
\caption{Angular power of the redshift anisotropy that an expansion tracking the matter distribution would produce, from \eqref{eq:limber} on the Planck 2018 linear matter spectrum, for three smoothing radii (red), against the temperature power spectrum Planck measures (blue)~\citep{Planck2018V}. Both are in units of $(\Delta T/T)^{2}$. The prediction exceeds the measurement at every multipole for every radius shown: by a factor of at least $2\times10^{4}$ for $R=4$ and $8\,h^{-1}$Mpc, and of at least $800$ for $R=20\,h^{-1}$Mpc, where the smoothing begins to remove the smallest angular scales.}
\label{fig:floor}
\end{figure}

The choice of $R$ does not rescue the hypothesis. The radius is the scale on which the rate is taken to follow the density, and the rms falls slowly as it grows: $1.9\times10^{-2}$ at $R=4$, $1.1\times10^{-2}$ at $20$, $6.0\times10^{-3}$ at $50$, $3.3\times10^{-3}$ at $100$ and $9.6\times10^{-4}$ at $300\,h^{-1}$Mpc---still $24$ times the whole measured anisotropy at a smoothing scale far above any on which the matter is inhomogeneous. As the radius grows, the smoothing removes the small angular scales and the excess moves to the large ones: at $R=50$ and $100\,h^{-1}$Mpc the predicted power still exceeds the measured at every multipole to $\ell=2500$, and at $R=300\,h^{-1}$Mpc it exceeds it at every multipole up to $\ell\approx490$, by a factor of more than $60$ up to $\ell=100$.

\subsection{Why the number is a floor}

Each approximation in the estimate lowers it:
\begin{itemize}
\item \emph{Linear theory.} The matter contrast on the scales that dominate \eqref{eq:sigmapath} is nonlinear at low redshift, where nonlinear growth raises the small-scale power above the linear spectrum used here; the void and wall contrasts the backreaction programmes invoke are of order unity.
\item \emph{The local rate.} Equation~\eqref{eq:localrate} is the linear limit of the separate-universe rate; the exact rate of a nonlinear underdensity departs from the mean by more, not less.
\item \emph{The comparison.} The measured anisotropy includes the primary acoustic signal, which the hypothesis does not remove; the redshift scatter it predicts must fit inside the measured power together with everything else on the sky, so comparing it with the whole measured spectrum, as though the primordial and Sachs--Wolfe contributions were absent, understates the exclusion.
\end{itemize}
The value~\eqref{eq:value} carries the correlations of the matter field along the ray, which the measured spectrum shows and structure formation in general relativity requires. The most accommodating value drops them: light crossing independent cells of uniform density, each of length $R$ and with its own contrast $\sigma_R(z)$, as many of them as the finite distance to last scattering holds---about $1200$ at $R=8\,h^{-1}$Mpc---so that the scatter averages down as far as the observable universe allows. Replacing the correlation length along the ray in~\eqref{eq:sigmapath} by $R$ gives $\sigma^{2}_{\mathrm{path}}=\int_0^{\chi_*}d\chi\,g^{2}(\chi)\,R\,\sigma_R^{2}(z)$, and $\sigma_{\mathrm{path}}\simeq8.9\times10^{-3}$, still some $220$ times the whole measured anisotropy. That value is unrealistically uniform, since the matter is not uncorrelated: correlation along the ray leaves fewer independent cells across the same finite path, and raises the scatter to~\eqref{eq:value}. The expectation general relativity itself supplies for the matter field thus moves the prediction further from the measured sky, not toward it. The Limber approximation in \eqref{eq:limber} is inaccurate at the lowest multipoles, by tens of per cent at $\ell\lesssim10$, which is immaterial against margins of four orders of magnitude; \eqref{eq:sigmapath} rests on the same narrow-correlation approximation along the ray, which holds because the correlation length of the smoothed field is short against the distances over which $g$ and $P$ change. The comparison is a lower bound against a measurement, which is what makes it robust: the exclusion cannot be weakened by a different estimate, only by a different physical claim.

\subsection{Distinct from the Sachs--Wolfe effect}\label{sec:sw}

It is essential that this is a distinct effect from the Sachs--Wolfe and integrated Sachs--Wolfe anisotropies~\citep{Sachs1967, Rees1968}, not a re-derivation of them. When the matter's lumpiness clusters upon a single background $a(t)$, it is carried by a metric potential $\Phi$. In \eqref{eq:redshiftrate} the local expansion of the matter still varies with its density, but the shear term no longer vanishes, and to linear order the two combine into a total derivative along the ray: the redshift a photon gains descending into a potential well is returned climbing out, the contributions \emph{telescope}, and what remains is the endpoint potential ($\sim\!10^{-5}$) and the small integrated term. The floor computed here is the contrary case---a genuine differential \emph{expansion}, a cosmological redshift that accumulates because there is nothing to cancel against. The observed anisotropy thus selects the telescoping, uniform-expansion picture over the accumulating, differential one. The isotropy of the CMB is not accounted for by large-scale statistical homogeneity~\citep{Scrimgeour2012}---that homogeneity is fully consistent with the differential-expansion hypothesis, which would still produce the scatter of~\eqref{eq:value}---but by uniform expansion in a cosmic time.

It is the scatter, not the mean, that is at issue. The redshift along a long ray may track the volume-averaged expansion~\citep{Rasanen2009}; the floor is the fluctuation about that average, computed. Under a matter-tracking rate it is of order $10^{-2}$ along each ray, and it is this direction-to-direction scatter, not the mean redshift, that the measured isotropy excludes.

A differential-expansion programme can evade the floor only by requiring the local rates to average to a common value in every direction: a uniform large-scale Hubble flow. But that requirement is itself a global uniform expansion in a cosmic time, the very structure such programmes set out to dispense with; imposing it to recover the observed isotropy reinstates the cosmic rest frame and the cosmic time it carries. Either way, the observed isotropy of the CMB redshift forces a uniform cosmic expansion.

\subsection{Timescape}\label{sec:timescape}

The timescape model of Section~\ref{sec:backreaction} falls inside the hypothesis, its void and wall rates differing by a third. Its postulate that each region measures the same rate in its own proper time does not remove the differential expansion from the redshift. By~\eqref{eq:redshiftrate} a photon crossing a region gains the local rate per unit of local proper time, and per unit of volume-average time it spends less local proper time crossing a wall than a void; the redshift it gains per unit of volume-average time is therefore the region's own rate, $H_{\rm v}$ or $H_{\rm w}$, while the clock rates enter only at the ends of the path, through the acceleration term. The mean of that accumulation is exactly the model's own redshift relation~\citep[eq.~38]{Wiltshire2009}, and its scatter is of the floor's kind: for void and wall cells drawn independently along each ray, the most accommodating case, the rms of $\ln(1+z)$ across directions is $2.4\times10^{-2}$, $4.2\times10^{-2}$ and $5.4\times10^{-2}$ for cells of $10$, $30$ and $50\,h^{-1}$\,Mpc, some $600$ to $1300$ times the whole measured anisotropy. Its variance is the backreaction $\mathcal{Q}$ integrated along the ray, weighted by the time light takes to cross each cell: the term that gives the model its acceleration is the term that scatters its redshifts. The timescape papers compute observables on the average geometry~\citep{Wiltshire2009, NazerWiltshire2015, Seifert2025} and do not compute this scatter. Uniformity in local proper time is not uniformity along the journey of the light, and it is the latter that the isotropy of the CMB measures.

\section{What the exclusion establishes}\label{sec:establishes}

The floor closes the statistical escape, and one further escape is logically available and must be closed separately. A finely tuned, spherically symmetric inhomogeneity centred on the observer would produce isotropic redshift about us with no global uniform expansion; being structured rather than statistical, it is not touched by the floor. It is closed instead by the Copernican principle together with the independent isotropy of the \emph{expansion history} itself: standard rulers (baryon acoustic oscillations~\citep{Eisenstein2005}) and standard candles (type~Ia supernovae~\citep{Riess1998, Perlmutter1999}) reconstruct the expansion along many directions and agree across the sky, so the uniformity is not the artefact of a privileged centre~\citep{Barrett2000, Clarkson2010, Maartens2011}. Gpc-scale voids centred on us, the form this escape takes in the literature, are excluded directly by the kinetic Sunyaev--Zel'dovich effect and by the combined distance data~\citep{ZhangStebbins2011, Moss2011}. With the statistical alternative excluded by the floor and the centred one by the Copernican step, the disjunction is exhausted: the observed isotropy forces uniform cosmic expansion.

What that exclusion reaches is the point of the paper. The region the data reject is not a parameter corner but the entire class of histories lacking a single global scale factor on a single global time. To reject that class is to select a global, time-ordered expansion, which is the cosmic foliation made dynamical. And by Section~\ref{sec:order} the foliation is logically prior to everything else the standard model assumes---``space'' is one of its slices, ``expansion'' an ordered family of them, ``isotropy'' and ``homogeneity'' properties of a slice---none of which is even statable until the foliation is in hand. A single datum therefore establishes, from the bottom up, the cosmic time on which the rest stands, the uniformity of the expansion on its slices, and---within the observable region and under the Copernican principle---the maximal symmetry of those slices. The model's foundational kinematics are not idealizations the data tolerate; they are what the data force, the one coherent alternative excluded. Unlike the stability arguments of Section~\ref{sec:tension}, the floor does not assume a Friedmann structure in its input; it derives the uniformity of the expansion from the redshift isotropy directly, so the uniform expansion that one side of the backreaction debate assumes and the other questions is fixed by measurement whichever way the averaging argument is settled.

The scope is worth marking. This contradicts none of the standard model's fitted predictions; with its assumptions in place it reproduces the observed expansion history, which is why it stands. What the measurement changes is the standing of those assumptions: the cosmic time and the uniformity of its expansion are forced from the bottom up, prior to and independently of the dynamical fit. The temporal claim---uniform expansion throughout the history probed---is clean, being the integral from decoupling to the present along the observed rays; the extension to a globally maximally symmetric space rests on the Copernican principle and extrapolates beyond the directly constrained region. And what the isotropy measures is the uniformity of the expansion rate along the observed rays; it does not measure, and this paper does not claim, anything about which geometry carries that expansion or how the matter within it is arranged.

\section{What stands and what falls}\label{sec:stands}

The floor constrains the expansion \emph{rate}, not the matter content. Space can be locally curved and inhomogeneous, with the matter density varying from place to place on each slice, and that lumpiness is untouched by the argument. What is excluded is an inhomogeneous expansion rate \emph{tracking} the lumpiness. Uniform expansion and lumpy matter are consistent.

The measurement leaves two things standing. The curvature stands in full: the lumpiness curves spacetime, and that curvature is the potential $\Phi$ of Section~\ref{sec:sw}, whose effect on the redshift telescopes along each ray. The Friedmann equation~\eqref{eq:friedmann} stands as a relation between averages, the mean density and the rate on a slicing already given. What does not stand is the step between them, the local form of Section~\ref{sec:twoforms}: that the density sets the rate region by region. The Raychaudhuri equation~\eqref{eq:raychaudhuri} is not what fails; it holds in every solution of general relativity. What fails is the claim that the rate light records is the matter's local expansion. The standard model's own perturbation theory is consistent with the floor for exactly this reason: the matter's local expansion varies with its density, but the shear that general relativity pairs with it cancels the variation out of the redshift, so the rate the light records is the uniform rate at which the slices of cosmic time advance. The curvature of the matter rides on that advance instead of setting it.

What does set the rate is then the global form of the principle. In the Friedmann models it is the average density on the slice, however lumpy the matter is on it, and the cosmic time and slices that form presupposes now rest on measurement rather than on posit. One technical qualification belongs here. The Friedmann models add a posit beyond the cosmic time itself: that the slices of constant cosmic time are orthogonal to the matter's worldlines, so that cosmic time is the proper time of the matter. Strictly, the redshift relation the measurement constrains does not by itself require this. It is also satisfied by a special non-synchronous case, in which every comoving worldline crosses the slices of uniform expansion in the same way, the slices being tilted from the matter's rest frame; an example whose scale factor is identical to that of flat $\Lambda$CDM is given by \citet{Janzen2015}. That case carries a caveat. On the metric of the example the comoving congruence is sheared, and null geodesics of that metric would return an anisotropic redshift; the case reproduces the observed isotropy only if the light is assigned to a different congruence, the worldlines at rest in the slices, by a reassignment of the metric's causal structure. Were it realized, events on one slice of constant cosmic time would be asynchronous in the matter's rest frame. We record the case for technical completeness only; nothing in this paper depends on it or argues for it.

The result saddles a theory without preferred frames with a preferred time, which may seem a high price to pay against the theory's elegance. Feynman's rule settles the order of priority: however beautiful a theory is, if it disagrees with experiment it is wrong~\citep{Feynman1965}. Here the disagreement does not fall on general relativity, whose equations admit the foliation the data select. It falls on two expectations about it: that no foliation would be distinguished, and that the density would set the expansion where it is.

\section{The measured foliation and what follows}\label{sec:follows}

\subsection{The order of construction reversed}

The construction of Section~\ref{sec:order} posited a cosmic time, its slices, and their isotropy and homogeneity so that the dynamical relation~\eqref{eq:friedmann} could be stated. The measurement reverses that order. The cosmic time and the uniformity of its expansion are measured, and with the Copernican principle so are the isotropy and homogeneity of its slices, while the local form of the dynamical principle, the one form that could have been tested without them, is what fails.

\subsection{The foliation in general relativity}

General relativity is generally covariant. A spacetime can be sliced into spacelike hypersurfaces in many ways, and the theory does not by itself make any slicing physically special. In the $3{+}1$ form of the theory~\citep{ADM1962} a slicing is described by two functions: the lapse $N$, the proper time elapsed along the normals between neighbouring slices, and the shift $N^{i}$, the motion of the spatial coordinates from one slice to the next relative to those normals. In the canonical theory both are Lagrange multipliers, freely specifiable. Numerical relativity chooses them for stability, with maximal or 1+log slicing~\citep{BaumgarteShapiro2010}, and the absence of any preferred foliation is the root of the problem of time in canonical quantum gravity~\citep{Kuchar1992, Isham1993}. Programmes that do posit a preferred foliation posit it on theoretical grounds: the constant-mean-curvature slicing that defines York's extrinsic time~\citep{York1972}, which shape dynamics takes as fundamental~\citep{Gomes2011}; the preferred foliation of Ho\v{r}ava--Lifshitz gravity~\citep{Horava2009}; and the dynamical preferred frame of Einstein--aether theory~\citep{Jacobson2001}.

In these terms, our universe has a cosmic time whose slices carry a uniform expansion: the redshift between any two comoving worldlines is the ratio of one scale factor at the cosmic times of reception and emission, whatever the direction, to the accuracy of the floor along our past light cone and, with the Copernican step of Section~\ref{sec:establishes}, across each slice. The slices are the level sets of that scale factor, and so of the temperature of the CMB in its own rest frame, which is a scalar: every observer at an event recovers the same value from the angular pattern they measure, whatever their motion. The slicing is therefore invariant rather than a convention of any observer's frame. Where the matter's worldlines are orthogonal to the slices, as the Friedmann models posit, the expansion of the matter is the mean curvature of the slices, so slices of uniform expansion are slices of constant mean curvature: the foliation the data select is York's. How the rate of comoving clocks varies against cosmic time is a separate matter. Across the last-scattering surface the measured anisotropy bounds it at the $10^{-5}$ level of the gravitational potentials; at the observer's end it is common to every direction, so isotropy seen from one place does not by itself fix how clock rates in different regions compare today. The slicing is not chosen. It is singled out by the content of the universe and located by measurement.

\subsection{What is unchanged}

All of special- and general-relativistic kinematics holds on this structure. A clock deep in a potential well runs slow by the gravitational redshift. A clock moving through the cosmic rest frame runs slow by its Lorentz factor. Observers moving relative to that frame who synchronize their clocks by light signals obtain slices tilted from the frame's own, by $vx/c^{2}$ to first order, and described on the cosmic slicing their frames carry a nonzero shift. Proper times, and the times read by relatively moving synchronized clocks, track differently, exactly as relativity requires, within a structure whose common time is measured. Local Lorentz invariance is untouched: no local experiment picks out the cosmic frame, which is identified only by the large-scale content of the universe.

\subsection{Simultaneity}

The relativity of simultaneity is a statement about synchrony---which distant clocks, set by a given convention, read the same---and it holds. It does not follow that a universe whose content distinguishes a slicing has no physically distinguished relation of simultaneity, and the data show that ours distinguishes one: the slices of constant cosmic time, on which the expansion is uniform. This bears directly on the arguments that draw a metaphysical conclusion from the relativity of simultaneity. Rietdijk and Putnam argued that if the simultaneity of every inertial frame is equally valid, events on all of them are equally real, so that the future is as determinate as the past~\citep{Rietdijk1966, Putnam1967}; Penrose's version sets two people passing in a street whose planes of simultaneity reach events in the Andromeda galaxy days apart~\citep{Penrose1989}. The premise is that no frame is physically distinguished, and Stein showed early that the relativistic structure alone does not license the conclusion~\citep{Stein1968}. What the data add is that in our universe the premise is false: the cosmic foliation is physically distinguished, by the content of the universe, and located by measurement. This does not settle the metaphysics of time; it removes the physical premise these arguments rest on. G\"odel objected that a cosmic time defined by the matter depends on how the matter happens to be arranged~\citep{Godel1949}; the measurement makes that dependence an empirical matter, settled for our universe. Accounts of an objective present built on cosmic time are developed by \citet{Ellis2006}, \citet{EllisRothman2010} and the contributors to \citet{CraigSmith2007}. Synchrony is relative; simultaneity on the cosmic foliation is absolute, and measured. The measurement does not settle which four-geometry carries that foliation, or how matter sits on it, and these are not claimed here.

\subsection{The cosmic rest frame}

The rest frame of the CMB is located by the same radiation: our own motion through it is read directly from the dipole, $(369.82\pm0.11)$\,km/s toward Crater~\citep{Planck2018I}. The dipole was predicted as the signature of that motion~\citep{PeeblesWilkinson1968}, detected~\citep{Smoot1977} and is now measured to a fraction of a per cent, and the aberration and Doppler boosting it imprints on the small-scale anisotropies confirm it as a motion~\citep{Planck2013XXVII}. Whether the mean rest frame of the matter coincides with that of the CMB is tested independently by the dipole in source counts~\citep{EllisBaldwin1984}, where an amplitude exceeding the kinematic expectation has been reported~\citep{Secrest2021, Secrest2025}. That test bears on how the matter's rest frame stands to the cosmic foliation, the question of the technical qualification in Section~\ref{sec:stands}, and not on the uniformity of the expansion measured here.

\subsection{The history closed}

The measurement closes the arc of Section~\ref{sec:history}. What Einstein posited as an approximation on stellar velocities, what de~Sitter placed beyond physical argument, and what Eddington defended in the declared absence of evidence, is what the redshift isotropy now forces. The assumption that ran against the spirit of relativity was correct, and the objection that it gave time ``a separate position'' was correct too---as a description of the structure, though not as a reason to reject it. The principle Friedmann drew from Einstein's model survives in its global form, now standing on a measured foundation, and fails in the local form that general relativity's locality seemed to promise. The parallax of this paper's title carries the same lesson. Stellar parallax went undetected for two millennia and was read as telling against the Earth's motion; an absent discriminator is not an absent fact, and whether a discriminator is even available is a question about the picture in hand rather than about the world.

\section{Conclusion}\label{sec:conclusion}

Relativistic cosmology was built on the principle that the matter determines how the universe expands. In its global form the principle is stated on a cosmic time and uniform slices that must be posited first; in its local form it presupposes nothing, and predicts that the expansion light records follows the matter it crosses. The isotropy of the CMB is a measurement of the expansion the radiation crossed, not only of the surface that emitted it. Read that way, it excludes the local form: an expansion tracking the matter would imprint redshift anisotropy at the $10^{-2}$ level, exceeding the measured temperature power at every multipole Planck observes by more than four orders of magnitude at $8\,h^{-1}$Mpc smoothing and remaining excluded for any smoothing scale up to $300\,h^{-1}$Mpc. The timescape model, through its own equations, exceeds it by several hundred to a thousand times. In our universe the local matter density does not set the rate of cosmic expansion. With the observer-centred alternative closed by the Copernican principle and the isotropy of the expansion history, the data select a single global expansion in a single cosmic time, and the structure the standard model posits before its dynamics is measured. As stellar parallax once converted the Earth's motion from an interpretive option into a measured fact, the isotropy of the cosmological redshift converts the cosmic foliation from a modelling convenience into a measured feature of the world.

\begin{acknowledgments}
Computations used \textsc{camb} and the Planck 2018 data products, based on observations obtained with Planck (\url{http://www.esa.int/Planck}), an ESA science mission with instruments and contributions directly funded by ESA Member States, NASA, and Canada. Large language model assistance (Anthropic's Claude) was used in preparing this manuscript.
\end{acknowledgments}

\facilities{Planck}

\software{CAMB \citep{Lewis2000}, NumPy, SciPy, Matplotlib}

\begin{thebibliography}{}

\bibitem[{Arnowitt} {et~al.}(1962)]{ADM1962} Arnowitt, R., Deser, S., \& Misner, C.~W. 1962, in Gravitation: An Introduction to Current Research, ed. L.~Witten (New York: Wiley), 227, arXiv:gr-qc/0405109
\bibitem[{Barrett} \& {Clarkson}(2000)]{Barrett2000} Barrett, R.~K., \& Clarkson, C.~A. 2000, CQGra, 17, 5047, doi:\doi{10.1088/0264-9381/17/23/315}
\bibitem[{Baumgarte} \& {Shapiro}(2010)]{BaumgarteShapiro2010} Baumgarte, T.~W., \& Shapiro, S.~L. 2010, Numerical Relativity: Solving Einstein's Equations on the Computer (Cambridge: Cambridge Univ. Press)
\bibitem[{Bessel}(1838)]{Bessel1838} Bessel, F.~W. 1838, AN, 16, 65, doi:\doi{10.1002/asna.18390160502}
\bibitem[{Buchert}(2000)]{Buchert2000} Buchert, T. 2000, GReGr, 32, 105, doi:\doi{10.1023/A:1001800617177}
\bibitem[{Buchert}(2008)]{Buchert2008} Buchert, T. 2008, GReGr, 40, 467, doi:\doi{10.1007/s10714-007-0554-8}
\bibitem[{Buchert} \& {R{\"a}s{\"a}nen}(2012)]{BuchertRasanen2012} Buchert, T., \& R{\"a}s{\"a}nen, S. 2012, ARNPS, 62, 57, doi:\doi{10.1146/annurev.nucl.012809.104435}
\bibitem[{Buchert} {et~al.}(2015)]{Buchert2015} Buchert, T., Carfora, M., Ellis, G.~F.~R., {et~al.} 2015, CQGra, 32, 215021, doi:\doi{10.1088/0264-9381/32/21/215021}
\bibitem[{Clarkson} {et~al.}(2008)]{Clarkson2010} Clarkson, C., Bassett, B., \& Lu, T.~H.-C. 2008, PhRvL, 101, 011301, doi:\doi{10.1103/PhysRevLett.101.011301}
\bibitem[{Clarkson} {et~al.}(2011)]{Clarkson2011} Clarkson, C., Ellis, G., Larena, J., \& Umeh, O. 2011, RPPh, 74, 112901, doi:\doi{10.1088/0034-4885/74/11/112901}
\bibitem[{Collins} \& {Hawking}(1973)]{CollinsHawking1973} Collins, C.~B., \& Hawking, S.~W. 1973, MNRAS, 162, 307, doi:\doi{10.1093/mnras/162.4.307}
\bibitem[{Craig} \& {Smith}(2007)]{CraigSmith2007} Craig, W.~L., \& Smith, Q., eds. 2007, Einstein, Relativity and Absolute Simultaneity (London: Routledge), doi:\doi{10.4324/9780203700051}
\bibitem[{de Sitter}(1917)]{deSitter1917} de Sitter, W. 1917, KNAB, 19, 1217
\bibitem[{Eddington}(1920)]{Eddington1920} Eddington, A.~S. 1920, Space, Time and Gravitation (Cambridge: Cambridge Univ. Press)
\bibitem[{Eddington}(1930)]{Eddington1930} Eddington, A.~S. 1930, MNRAS, 90, 668, doi:\doi{10.1093/mnras/90.7.668}
\bibitem[{Ehlers} {et~al.}(1968)]{EGS1968} Ehlers, J., Geren, P., \& Sachs, R.~K. 1968, JMP, 9, 1344, doi:\doi{10.1063/1.1664720}
\bibitem[{Einstein}(1917)]{Einstein1917} Einstein, A. 1917, Sitzungsber. K\"onigl. Preuss. Akad. Wiss. (Berlin), 142
\bibitem[{Eisenstein} {et~al.}(2005)]{Eisenstein2005} Eisenstein, D.~J., Zehavi, I., Hogg, D.~W., {et~al.} 2005, ApJ, 633, 560, doi:\doi{10.1086/466512}
\bibitem[{Ellis}(1971)]{Ellis1971} Ellis, G.~F.~R. 1971, in Proc. Int. School of Physics ``Enrico Fermi'', Course 47: General Relativity and Cosmology, ed. R.~K.~Sachs (New York: Academic Press), 104; reprinted 2009, GReGr, 41, 581
\bibitem[{Ellis}(1984)]{Ellis1984} Ellis, G.~F.~R. 1984, in General Relativity and Gravitation, ed. B.~Bertotti, F.~de Felice, \& A.~Pascolini (Dordrecht: Reidel), 215, doi:\doi{10.1007/978-94-009-6469-3_14}
\bibitem[{Ellis}(2006)]{Ellis2006} Ellis, G.~F.~R. 2006, GReGr, 38, 1797, doi:\doi{10.1007/s10714-006-0332-z}
\bibitem[{Ellis} \& {Baldwin}(1984)]{EllisBaldwin1984} Ellis, G.~F.~R., \& Baldwin, J.~E. 1984, MNRAS, 206, 377, doi:\doi{10.1093/mnras/206.2.377}
\bibitem[{Ellis} \& {Rothman}(2010)]{EllisRothman2010} Ellis, G.~F.~R., \& Rothman, T. 2010, IJTP, 49, 988, doi:\doi{10.1007/s10773-010-0278-5}
\bibitem[{Ellis} {et~al.}(2012)]{EllisMaartensMacCallum2012} Ellis, G.~F.~R., Maartens, R., \& MacCallum, M.~A.~H. 2012, Relativistic Cosmology (Cambridge: Cambridge Univ. Press)
\bibitem[{Feynman}(1965)]{Feynman1965} Feynman, R.~P. 1965, The Character of Physical Law (Cambridge, MA: MIT Press)
\bibitem[{Friedmann}(1922)]{Friedmann1922} Friedmann, A. 1922, ZPhy, 10, 377, doi:\doi{10.1007/BF01332580}
\bibitem[{Friedmann}(1924)]{Friedmann1924} Friedmann, A. 1924, ZPhy, 21, 326, doi:\doi{10.1007/BF01328280}
\bibitem[{G{\"o}del}(1949)]{Godel1949} G{\"o}del, K. 1949, in Albert Einstein: Philosopher-Scientist, ed. P.~A.~Schilpp (Evanston: Library of Living Philosophers), 555
\bibitem[{Gomes} {et~al.}(2011)]{Gomes2011} Gomes, H., Gryb, S., \& Koslowski, T. 2011, CQGra, 28, 045005, doi:\doi{10.1088/0264-9381/28/4/045005}
\bibitem[{Green} \& {Wald}(2011)]{GreenWald2011} Green, S.~R., \& Wald, R.~M. 2011, PhRvD, 83, 084020, doi:\doi{10.1103/PhysRevD.83.084020}
\bibitem[{Green} \& {Wald}(2014)]{GreenWald2014} Green, S.~R., \& Wald, R.~M. 2014, CQGra, 31, 234003, doi:\doi{10.1088/0264-9381/31/23/234003}
\bibitem[{Ho{\v r}ava}(2009)]{Horava2009} Ho{\v r}ava, P. 2009, PhRvD, 79, 084008, doi:\doi{10.1103/PhysRevD.79.084008}
\bibitem[{Hubble}(1929)]{Hubble1929} Hubble, E. 1929, PNAS, 15, 168, doi:\doi{10.1073/pnas.15.3.168}
\bibitem[{Isham}(1993)]{Isham1993} Isham, C.~J. 1993, in Integrable Systems, Quantum Groups, and Quantum Field Theories, ed. L.~A.~Ibort \& M.~A.~Rodr\'iguez (Dordrecht: Kluwer), 157, arXiv:gr-qc/9210011
\bibitem[{Jacobson} \& {Mattingly}(2001)]{Jacobson2001} Jacobson, T., \& Mattingly, D. 2001, PhRvD, 64, 024028, doi:\doi{10.1103/PhysRevD.64.024028}
\bibitem[{Janzen}(2014)]{JanzenEinsteinCosmology} Janzen, D. 2014, arXiv:1402.3212
\bibitem[{Janzen}(2015)]{Janzen2015} Janzen, D. 2015, in Questioning the Foundations of Physics, ed. A.~Aguirre, B.~Foster, \& Z.~Merali (Cham: Springer), 103, arXiv:1301.5980
\bibitem[{Kolb} {et~al.}(2006)]{Kolb2006} Kolb, E.~W., Matarrese, S., \& Riotto, A. 2006, NJPh, 8, 322, doi:\doi{10.1088/1367-2630/8/12/322}
\bibitem[{Kucha{\v r}}(1992)]{Kuchar1992} Kucha{\v r}, K.~V. 1992, in Proc. 4th Canadian Conf. on General Relativity and Relativistic Astrophysics, ed. G.~Kunstatter, D.~E.~Vincent, \& J.~G.~Williams (Singapore: World Scientific), 211
\bibitem[{Lema{\^i}tre}(1927)]{Lemaitre1927} Lema{\^i}tre, G. 1927, ASSB, 47, 49
\bibitem[{Lewis} {et~al.}(2000)]{Lewis2000} Lewis, A., Challinor, A., \& Lasenby, A. 2000, ApJ, 538, 473, doi:\doi{10.1086/309179}
\bibitem[{Limber}(1953)]{Limber1953} Limber, D.~N. 1953, ApJ, 117, 134, doi:\doi{10.1086/145672}
\bibitem[{Maartens}(2011)]{Maartens2011} Maartens, R. 2011, RSPTA, 369, 5115, doi:\doi{10.1098/rsta.2011.0289}
\bibitem[{Moss} {et~al.}(2011)]{Moss2011} Moss, A., Zibin, J.~P., \& Scott, D. 2011, PhRvD, 83, 103515, doi:\doi{10.1103/PhysRevD.83.103515}
\bibitem[{Nazer} \& {Wiltshire}(2015)]{NazerWiltshire2015} Nazer, M.~A., \& Wiltshire, D.~L. 2015, PhRvD, 91, 063519, doi:\doi{10.1103/PhysRevD.91.063519}
\bibitem[{Nilsson} {et~al.}(1999)]{Nilsson1999} Nilsson, U.~S., Uggla, C., Wainwright, J., \& Lim, W.~C. 1999, ApJL, 522, L1, doi:\doi{10.1086/312209}
\bibitem[{North}(1965)]{north1965} North, J.~D. 1965, The Measure of the Universe: A History of Modern Cosmology (Oxford: Oxford Univ. Press)
\bibitem[{Peebles}(1980)]{Peebles1980} Peebles, P.~J.~E. 1980, The Large-Scale Structure of the Universe (Princeton: Princeton Univ. Press)
\bibitem[{Peebles} \& {Wilkinson}(1968)]{PeeblesWilkinson1968} Peebles, P.~J.~E., \& Wilkinson, D.~T. 1968, PhRv, 174, 2168, doi:\doi{10.1103/PhysRev.174.2168}
\bibitem[{Penrose}(1989)]{Penrose1989} Penrose, R. 1989, The Emperor's New Mind (Oxford: Oxford Univ. Press)
\bibitem[{Perlmutter} {et~al.}(1999)]{Perlmutter1999} Perlmutter, S., Aldering, G., Goldhaber, G., {et~al.} 1999, ApJ, 517, 565, doi:\doi{10.1086/307221}
\bibitem[{Planck Collaboration} {et~al.}(2014)]{Planck2013XXVII} {Planck Collaboration}, Aghanim, N., Armitage-Caplan, C., {et~al.} 2014, A\&A, 571, A27, doi:\doi{10.1051/0004-6361/201321556}
\bibitem[{Planck Collaboration} {et~al.}(2020a)]{Planck2018I} {Planck Collaboration}, Aghanim, N., Akrami, Y., {et~al.} 2020a, A\&A, 641, A1, doi:\doi{10.1051/0004-6361/201833880}
\bibitem[{Planck Collaboration} {et~al.}(2020b)]{Planck2018V} {Planck Collaboration}, Aghanim, N., Akrami, Y., {et~al.} 2020b, A\&A, 641, A5, doi:\doi{10.1051/0004-6361/201936386}
\bibitem[{Planck Collaboration} {et~al.}(2020c)]{Planck2018VI} {Planck Collaboration}, Aghanim, N., Akrami, Y., {et~al.} 2020c, A\&A, 641, A6, doi:\doi{10.1051/0004-6361/201833910}
\bibitem[{Putnam}(1967)]{Putnam1967} Putnam, H. 1967, J. Philos., 64, 240, doi:\doi{10.2307/2024493}
\bibitem[{R{\"a}s{\"a}nen}(2004)]{Rasanen2004} R{\"a}s{\"a}nen, S. 2004, JCAP, 2004, 003, doi:\doi{10.1088/1475-7516/2004/02/003}
\bibitem[{R{\"a}s{\"a}nen}(2009)]{Rasanen2009} R{\"a}s{\"a}nen, S. 2009, JCAP, 2009, 011, doi:\doi{10.1088/1475-7516/2009/02/011}
\bibitem[{Raychaudhuri}(1955)]{Raychaudhuri1955} Raychaudhuri, A. 1955, PhRv, 98, 1123, doi:\doi{10.1103/PhysRev.98.1123}
\bibitem[{Rees} \& {Sciama}(1968)]{Rees1968} Rees, M.~J., \& Sciama, D.~W. 1968, Natur, 217, 511, doi:\doi{10.1038/217511a0}
\bibitem[{Riess} {et~al.}(1998)]{Riess1998} Riess, A.~G., Filippenko, A.~V., Challis, P., {et~al.} 1998, AJ, 116, 1009, doi:\doi{10.1086/300499}
\bibitem[{Rietdijk}(1966)]{Rietdijk1966} Rietdijk, C.~W. 1966, PhilSci, 33, 341, doi:\doi{10.1086/288106}
\bibitem[{Robertson}(1929)]{Robertson1929} Robertson, H.~P. 1929, PNAS, 15, 822, doi:\doi{10.1073/pnas.15.11.822}
\bibitem[{Robertson}(1933)]{Robertson1933} Robertson, H.~P. 1933, RvMP, 5, 62, doi:\doi{10.1103/RevModPhys.5.62}
\bibitem[{Robertson}(1935)]{Robertson1935} Robertson, H.~P. 1935, ApJ, 82, 284, doi:\doi{10.1086/143681}
\bibitem[{Saadeh} {et~al.}(2016)]{Saadeh2016} Saadeh, D., Feeney, S.~M., Pontzen, A., Peiris, H.~V., \& McEwen, J.~D. 2016, PhRvL, 117, 131302, doi:\doi{10.1103/PhysRevLett.117.131302}
\bibitem[{Sachs} \& {Wolfe}(1967)]{Sachs1967} Sachs, R.~K., \& Wolfe, A.~M. 1967, ApJ, 147, 73, doi:\doi{10.1086/148982}
\bibitem[{Scrimgeour} {et~al.}(2012)]{Scrimgeour2012} Scrimgeour, M.~I., Davis, T., Blake, C., {et~al.} 2012, MNRAS, 425, 116, doi:\doi{10.1111/j.1365-2966.2012.21402.x}
\bibitem[{Secrest} {et~al.}(2021)]{Secrest2021} Secrest, N.~J., von Hausegger, S., Rameez, M., {et~al.} 2021, ApJL, 908, L51, doi:\doi{10.3847/2041-8213/abdd40}
\bibitem[{Secrest} {et~al.}(2025)]{Secrest2025} Secrest, N.~J., von Hausegger, S., Rameez, M., Mohayaee, R., \& Sarkar, S. 2025, RvMP, 97, 041001, doi:\doi{10.1103/9ygx-z2yq}
\bibitem[{Seifert} {et~al.}(2025)]{Seifert2025} Seifert, A., Lane, Z.~G., Galoppo, M., Ridden-Harper, R., \& Wiltshire, D.~L. 2025, MNRAS, 537, L55, doi:\doi{10.1093/mnrasl/slae112}
\bibitem[{Smoot} {et~al.}(1977)]{Smoot1977} Smoot, G.~F., Gorenstein, M.~V., \& Muller, R.~A. 1977, PhRvL, 39, 898, doi:\doi{10.1103/PhysRevLett.39.898}
\bibitem[{Stein}(1968)]{Stein1968} Stein, H. 1968, J. Philos., 65, 5, doi:\doi{10.2307/2024512}
\bibitem[{Stoeger} {et~al.}(1995)]{Stoeger1995} Stoeger, W.~R., Maartens, R., \& Ellis, G.~F.~R. 1995, ApJ, 443, 1, doi:\doi{10.1086/175496}
\bibitem[{Walker}(1937)]{Walker1937} Walker, A.~G. 1937, Proc. London Math. Soc. (2), 42, 90
\bibitem[{Weyl}(1923)]{Weyl1923} Weyl, H. 1923, PhyZ, 24, 230; English translation 2009, GReGr, 41, 1661, doi:\doi{10.1007/s10714-009-0826-6}
\bibitem[{Wheeler} \& {Ford}(1998)]{Wheeler1998} Wheeler, J.~A., \& Ford, K. 1998, Geons, Black Holes, and Quantum Foam: A Life in Physics (New York: W.~W.~Norton), 235
\bibitem[{Wiltshire}(2007a)]{Wiltshire2007} Wiltshire, D.~L. 2007a, NJPh, 9, 377, doi:\doi{10.1088/1367-2630/9/10/377}
\bibitem[{Wiltshire}(2007b)]{Wiltshire2007PRL} Wiltshire, D.~L. 2007b, PhRvL, 99, 251101, doi:\doi{10.1103/PhysRevLett.99.251101}
\bibitem[{Wiltshire}(2009)]{Wiltshire2009} Wiltshire, D.~L. 2009, PhRvD, 80, 123512, doi:\doi{10.1103/PhysRevD.80.123512}
\bibitem[{York}(1972)]{York1972} York, J.~W. 1972, PhRvL, 28, 1082, doi:\doi{10.1103/PhysRevLett.28.1082}
\bibitem[{Zhang} \& {Stebbins}(2011)]{ZhangStebbins2011} Zhang, P., \& Stebbins, A. 2011, PhRvL, 107, 041301, doi:\doi{10.1103/PhysRevLett.107.041301}
\bibitem[{Zhou} {et~al.}(2025)]{Zhou2025} Zhou, A.~J., Dodelson, S., \& Scolnic, D. 2025, arXiv:2506.14878

\end{thebibliography}

\end{document}
